\chapter{Quick Start}
\label{ch:qtsl-quickstart}

In this chapter you will produce a machine-checkable certificate of a
quantitative claim about a stochastic system, in milliseconds, and see
exactly what any checker must re-derive from its bytes. You will meet
the three routes by which certificate bytes get checked in this
toolkit --- and where each route stands at this preview.
Then you will break a certificate from the published corpus and watch
an exact-arithmetic reader refuse it, with the violated obligation
named. The whole loop fits in one sitting; every command is
copy-pasteable, and every output shown is the verbatim transcript of
a real run --- regenerated by the book's own build, which executes
the printed commands and records the tool revisions and run date
beside each transcript in the source (one commodity arm64 box; your
timings will differ, your bytes will not).

\section{The claim}
\label{sec:qs-claim}

FrozenLake is Gymnasium's introductory gridworld: a $4\times4$ frozen
lake, sixteen states, four actions, a start square, four holes, one
goal. In the \emph{slippery} variant each action moves as intended
with probability $\tfrac13$ and slips to a perpendicular neighbour
otherwise, so every trajectory is a random walk biased by the policy,
and every reach probability is a rational number whose denominator is
a power of three.

The claim to certify: \emph{the recorded policy reaches the goal
within ten steps with probability at least}
\[
  p \;=\; \frac{815}{19683}
  \qquad\text{(exactly; } 19683 = 3^9\text{).}
\]
In QTSL's vocabulary this is a probabilistic deadline claim,
$\mathbb{P}_{\ge p}\!\left[\,\mathrm{F}_{\le 10}\;\mathit{goal}\,\right]$
--- the same claim species whose full campaign (deadline ladders,
dyadic rounding lanes, a family of environments) the Experiments
chapter (chapter~\ref{ch:experiments}) re-derives from committed
artefacts.
One disclosure before we start: $815/19683$ happens to be the exact
optimum of the ten-step problem, and the producer knows it (it runs an
independent floating-point sweep as a cross-instrument), but the
certificate below asserts only the lower bound. What a certificate
\emph{checks} is always narrower than what its producer \emph{knows};
keeping that line crisp is most of the craft, and the rest of this
book is about the vocabulary for widening what is checked.

\section{Setup}
\label{sec:qs-setup}

Three repositories, cloned side by side. \path{maitria-coalition/mtk}
is the producer toolkit and the only thing you build: a Rust
toolchain is its one hard prerequisite, and a GPU is optional
throughout (the accelerated lanes load CUDA dynamically and fall back
to the same exact host arithmetic --- same script, same answers).
\path{maitria-coalition/qtsl-experiments} is the reproduction spine
--- committed certificates, mutant batteries, receipts --- cloned for
its bytes; nothing to build. \path{maitria-coalition/qtsl} is the
repository this book ships in: the theories, the machine-checked
soundness developments, and the small reference instruments; nothing
to build there either --- the reader you will run in
section~\ref{sec:qs-break} is dependency-free TypeScript on bare Node
$\ge 22$.

\begin{Verbatim}[fontsize=\small,frame=leftline,rulecolor=\color{accentmid},framesep=3mm,xleftmargin=3mm]
git clone https://github.com/maitria-coalition/mtk
cargo build --release --manifest-path mtk/Cargo.toml
git clone https://github.com/maitria-coalition/qtsl-experiments
git clone https://github.com/maitria-coalition/qtsl
\end{Verbatim}

\section{Produce}
\label{sec:qs-produce}

\begin{Verbatim}[fontsize=\small,frame=leftline,rulecolor=\color{accentmid},framesep=3mm,xleftmargin=3mm]
$ mtk/target/release/gymdp produce --out out --only frozenlake-4x4
frozenlake-4x4: exact T=100 bound = 0.7441903 [0.005s]
receipts -> out/producer_receipts_rust.json (9 certs, 0 compared, 0 mismatched) [0.005s total]
\end{Verbatim}


Nine certificates in milliseconds: the exact deadline ladder
$T \in \{5, 10, 25, 50, 100\}$ and a ladder of dyadic-rounded
companions. (The \texttt{0.7441903} on the log line is a float
\emph{rendering} of the $T{=}100$ bound, printed for the human reading
the log; no float appears inside any certificate.) The document this
chapter cares about is the deadline-10 certificate,
\path{out/frozenlake-4x4/exact_T10.cbor}: 2{,}511 canonical bytes
carrying the claim, the transition kernel as 148 rows of exact
rationals, the recorded policy, and the eleven-by-sixteen value table
--- the complete witness data from which everything asserted can be
re-derived by something that has never met the producer. It is also,
byte for byte, the artefact committed in the experiments repository
(\texttt{sha256} prefix \texttt{66040e59}, either way --- compare
against \path{qtsl-experiments/gymnasium/out/frozenlake-4x4/}). In
this toolkit ``reproduces the published artefact'' is a statement
about equality of bytes, not about agreement to some tolerance ---
which is what lets a checker's verdict on the committed corpus and on
your fresh production be the same fact.

One vocabulary correction before anyone checks anything, because the
file extension just lied to you: this toolkit has no certificate file
format. In the shipped architecture a certificate is ordinary
database content --- rows of a signed, content-addressed geolog
commit, carried on the same rails as every other fact the system stores
(the Executive Summary's one-set-of-bytes, three-consumers stance
--- chapter~\ref{ch:exec-summary}; the Geolog part of this book is
that database's own manual). The theory defining those rows is in
this repository, as data: \path{theories/Qtsl.geolog}, with its
drift-alarmed battery beside it. The
sidecar file you just wrote is the producer's local carriage for that
content --- and the reproduction spine already commits this family
both ways: the sidecar artefacts, and the same certificates sealed as
rows of signed geolog commits in a \path{glc/} directory beside them,
with content identity between the two forms re-proven at every
production. What survives every carriage is the standard the content
answers to: a self-contained document of exact numbers, complete
enough to re-derive every obligation, with no side channel to
consult. You will read a certificate directly as rows, in plain text,
two sections from now.

\section{What a checker re-derives --- and who checks}
\label{sec:qs-check}

Any checker of this certificate, by any route, re-derives four
obligation families from the document alone, in exact rational
arithmetic:

\begin{itemize}
\item \emph{well-formedness} --- kernel probabilities lie in $[0,1]$
  and sum to one per state--action pair; indices are in range;
\item \emph{boundary} --- the value table's final row vanishes off the
  goal set;
\item \emph{step} --- each table entry is dominated by its expected
  successor value under the recorded policy's action,
  $V_t(s) \le \mathbb{E}\!\left[V_{t+1}\right]$;
\item \emph{claim} --- $p \le V_0(\mathit{init})$.
\end{itemize}

\noindent Chained, these force the table to be a genuine
under-approximation of the reach probability, so the claim follows
from the bytes alone. Nothing the producer computed is taken on
trust; the producer's speed lives entirely upstream of the
certificate. This division --- fast untrusted search, small
independent replay --- is the toolkit's organizing principle, and it
is why the producer is allowed to be a GPU-accelerated,
float-assisted, LLM-advised anything: a wrong producer costs a
rejected certificate, never a false verdict.

Who does the re-deriving is a deliberately plural answer: one set of
bytes, three consumers, each trading speed against the size of what
you must trust (the Executive Summary carries the full stance).
Their standing for \emph{this} certificate, dated to this preview:

\begin{itemize}
\item \emph{The geolog engine} --- the fast path: certificate rows
  chased, monitored, and enforced like any other rows;
  engineering-grade assurance, deliberately no mechanized soundness
  story. This certificate's row form --- one sealed commit carrying
  the system, the specification, the judgment, and the proof content
  --- is committed in the reproduction spine beside the sidecar
  artefacts, and the fast-path checker re-verifies all four
  obligation families from the sealed rows alone: one make target
  re-checks the family's containers against their manifest, expected
  verdicts included, and the Experiments chapter's timing tables
  carry the fast-path column from real runs.
\item \texttt{qtslab} --- the certified checker: a standalone binary
  whose end-to-end soundness theorem (machine-checked in HOL4,
  grounded through the CakeML verified compiler in the Arm
  instruction set's semantics) says that if it prints ACCEPT on a
  claim followed by an arbitrary stream of bytes, the claim is
  semantically true. Precisely: acceptance binds to validity of the
  accepted sequent under every admissible valuation of the hypernet
  the certificate itself declares (chapter~\ref{ch:qtsl-hypernets})
  --- the Bernstein Logic fragment,
  which is all its input format carries today; a certificate
  reaching beyond that fragment is refused, not silently accepted
  (the Executive Summary is explicit about the scope: the boundary
  is \texttt{qtslab}'s, not the row vocabulary's). A demonstrated
  route, assembling now: the same
  compiler-and-theorem recipe has already carried a complete
  predecessor checker to a machine-code soundness theorem on real
  hardware --- and that proven chain, re-run on its native ISA
  during this chapter's verification campaign, verified campaign
  fixtures in single-digit seconds each and refused all six mutants
  put to it.
\item \emph{The Lean route} --- the same bytes reconstructed as a
  theorem inside Lean's kernel, by reflection; expected to be too
  slow for machine-scale certificates, and kept in the roster to
  \emph{measure} that rather than assert it. It has been measured,
  twice, across two certificate-format generations: the Experiments
  chapter carries the scaling law, the admission gate fitted to it,
  and the expectation confirmed with its boundary stated.
\end{itemize}

\noindent No transcript in this section, on purpose: for the
certified and Lean routes, printing one would mean pretending they
are further along for this family than they are, and a book whose
subject is checkable evidence does not get to fake its own. The fast
path is runnable today --- the reproduction spine's make target
re-checks this family's committed containers, mutant battery
included, whenever you like. The checking you can run \emph{and read
whole} is the next section's, on a certificate family whose rows are
committed as plain text.

\section{Break one}
\label{sec:qs-break}

A checker that cannot say \emph{no} teaches nothing. The FrozenLake
family's sealed rows come with a committed mutant battery of their
own, but a sealed container is built for checkers' eyes; the
published-table campaign additionally stores its certificates as
plain-text row listings you can open and read --- so we break one of
those, and the lesson transfers intact.

That campaign (section~\ref{sec:exp-sweep} of
chapter~\ref{ch:experiments}) re-derived a peer-reviewed
neural-certificate results table in exact arithmetic, and
\texttt{faulty\_uniform} is its most instructive row: the one whose
exact two-sided bracket showed a published ``sound'' bound to be no
bound at all. Its upper-bound certificate is committed as rows in the
reference listing notation --- plain text you can open:

\begin{Verbatim}[fontsize=\small,frame=leftline,rulecolor=\color{accentmid},framesep=3mm,xleftmargin=3mm]
$ cd qtsl-experiments/supermartingales/out/sec3
$ sed -n '62p' certs/faulty_uniform/bound.rows.txt
op[  59] funa fun=G<stack/TConst> args=[q:16524448999/8406920828251] value=L58
\end{Verbatim}


\noindent One hundred and twelve rows like this one --- sorts, exact
simplex coordinates, Bernstein coefficients, and the proof structure
that cites them --- resolving their names against pinned theory
listings. Beside the corpus, the book's own repository carries a
reference reader and checker for this certificate family: three
TypeScript files, zero dependencies, readable in an afternoon
(\path{qtsl/checkers/supermartingales}). Point it at the certificate,
and at the campaign's committed mutant of the same certificate:

\begin{Verbatim}[fontsize=\small,frame=leftline,rulecolor=\color{accentmid},framesep=3mm,xleftmargin=3mm]
$ CHK=../../../../qtsl/checkers/supermartingales
$ node --experimental-strip-types $CHK/check_rows.ts \
    --stack $CHK/theories/BLstack.rows.txt --cert $CHK/theories/BLCert.rows.txt \
    certs/faulty_uniform/bound.rows.txt mutants/faulty_uniform/M1_false_bound.rows.txt
ACCEPT certs/faulty_uniform/bound.rows.txt proof_nodes=3
REJECT mutants/faulty_uniform/M1_false_bound.rows.txt berncone.coeff
\end{Verbatim}


\noindent Now break it yourself. The certificate survives a
coefficient perturbed in the sound direction --- slack is real, and
the checker is not grading effort --- so nudge the coefficient you
just looked at the \emph{other} way, by exactly
$1/8{,}406{,}920{,}828{,}251$:

\begin{Verbatim}[fontsize=\small,frame=leftline,rulecolor=\color{accentmid},framesep=3mm,xleftmargin=3mm]
$ sed 's|448999/|449000/|' certs/faulty_uniform/bound.rows.txt > my_broken.rows.txt
$ node --experimental-strip-types $CHK/check_rows.ts \
    --stack $CHK/theories/BLstack.rows.txt --cert $CHK/theories/BLCert.rows.txt \
    my_broken.rows.txt
REJECT my_broken.rows.txt berncone.coeff
\end{Verbatim}


Two things to notice. The refusal is \emph{witness-shaped}: the
identifier names the violated obligation family (membership of the
certificate polynomial's Bernstein coefficients in the claimed cone),
and the checker's failure record beneath it carries the offending
coefficient's index and exact value, because failure geometry is
routing advice for whoever must repair the certificate. And the
committed mutant battery is the checker's obligation list made
executable --- the battery, not the prose, gives the operational
statement of what an \texttt{ACCEPT} means (the Cookbook develops
this discipline as a pattern:
section~\ref{sec:ckbk-mutant-battery} of
chapter~\ref{ch:mtk-cookbook}).

One label, to keep the trust story crisp: the program you
just ran is a \emph{reference instrument} --- the reproduction
spine's own exactness oracle, differentially pinned against an
independent implementation across the campaign corpus and its mutant
battery --- and deliberately \emph{not} one of the three checking
routes. Its job is to make ``the battery must refuse'' an executable
fact, and to let you watch obligations re-derived from rows alone
through code small enough to read whole. When the stakes are real,
the same rows go to the routes of section~\ref{sec:qs-check}.

\section{The same loop from TypeScript}
\label{sec:qs-typescript}

The producer toolkit's TypeScript surface, \texttt{@maitria/mtk},
ships as the \path{ts/} package of the \path{maitria-coalition/mtk}
repository (napi bindings over the same Rust core --- there is one
implementation of every construction, and TypeScript reaches it
in-process). The shipped install route is the checkout you already
have from section~\ref{sec:qs-setup} --- it requires a Rust
toolchain, because the native addon builds from source:

\begin{Verbatim}[fontsize=\small,frame=leftline,rulecolor=\color{accentmid},framesep=3mm,xleftmargin=3mm]
cd mtk/ts && npm install && npm run build:native   # Rust toolchain required
cd ~/your-project && npm install ../mtk/ts         # npm links the package
\end{Verbatim}

\noindent (\texttt{@maitria/mtk} on the npm registry: arriving; the
checkout path above is the shipped route at this preview.) Node
$\ge 22.18$ throughout --- the package's declared \texttt{engines}
floor. The producing half of the loop again, this time
constructing the certificate yourself --- and building the mutant
with your own hands:

\begin{Verbatim}[fontsize=\small,frame=leftline,rulecolor=\color{accentmid},framesep=3mm,xleftmargin=3mm]
// quickstart.ts — produce the FrozenLake certificate, then break it.
import { writeFileSync } from "node:fs";
import { dp } from "@maitria/mtk";

// The environment: Gymnasium's FrozenLake-4x4, slippery (success rate 1/3).
const model = dp.frozenlake("4x4", true, "1/3");

// The claim: exact finite-horizon dynamic programming to deadline T = 10.
const res = model.exact(10);
console.log("claim:", res.bound);

// Emit the certificate: canonical bytes, encoded by the Rust core.
const meta = {
  envId: "FrozenLake-v1",
  mdpName: "frozenlake-4x4",
  repoSha: dp.gymSha(),
  kwargs: "",
};
const bytes = res.emit(meta);
writeFileSync("exact_T10.cbor", bytes);
console.log("wrote exact_T10.cbor:", bytes.length, "bytes");

// Now break it: raise one interior value-table entry to 1.
// (docMutate refuses unless the unmodified doc round-trips byte-identically.)
const broken = dp.docMutate(bytes, [
  { path: ["table", 1, 0, 0], setHex: "01" },
  { path: ["table", 1, 0, 1], setHex: "01" },
]);
writeFileSync("mutant_table_raised.cbor", broken);
console.log("wrote mutant_table_raised.cbor:", broken.length, "bytes");
\end{Verbatim}


\begin{Verbatim}[fontsize=\small,frame=leftline,rulecolor=\color{accentmid},framesep=3mm,xleftmargin=3mm]
$ npx tsx quickstart.ts
claim: 815/19683
wrote exact_T10.cbor: 2511 bytes
wrote mutant_table_raised.cbor: 2509 bytes
\end{Verbatim}


The bytes you just constructed are byte-identical to the ones
\texttt{gymdp} emitted in section~\ref{sec:qs-produce}, and therefore
to the committed corpus artefact --- the same \texttt{sha256}, all
three ways. And you now hold what the mutant-battery discipline calls
a matched pair: golden bytes, and a variant broken in precisely one
field (the mutation seam checks that the unmodified document
round-trips byte-identically before it will perform surgery, so your
mutant differs from the golden bytes exactly where you aimed). The
corpus commits this family's battery beside the artefacts, expected
verdicts in the receipts --- and the sealed row form carries a second
battery of thirteen row-form mutants, each one's expected verdict and
violated obligation named in its manifest. You already know from
section~\ref{sec:qs-break} what a refusal at the row level looks
like.

\section{What you just did, and where it goes}
\label{sec:qs-next}

You ran the trust architecture of this book once, in miniature: an
untrusted producer constructed complete witness data fast, and
reproduced a published artefact to the byte; you read the obligations
a checker must re-derive, and the current state of the three routes
that re-derive them; and on the published corpus you watched an
exact-arithmetic reader accept a certificate stored as rows and
refuse its mutants --- including yours --- with the violated
obligation named. Three widenings from here:

\begin{itemize}
\item \emph{Producing for real.} The Tutorial
  (chapter~\ref{ch:qtsl-tutorial}) teaches the producer's craft on the
  TypeScript surface: exactness as a type property, propose-then-dispose,
  specs and judgments, derivation rows from the Rule Catalogue, and
  refusal-driven iteration.
\item \emph{Richer claims.} Deadline reachability is one species. The
  Rule Catalogue (chapter~\ref{ch:qtsl-manual}) is the full inventory
  of certified inference; the fabrication-line worked example
  (chapter~\ref{ch:qtsl-minifab}) runs the stack over a model with
  batching, failures, and re-entrant flow.
\item \emph{Graded checking, priced.} The three routes of
  section~\ref{sec:qs-check} are not interchangeable, and the
  Experiments chapter's design keeps their price differences public:
  three wall-clocks per experiment --- produce, check-fast,
  check-certified --- each measured on its route's native substrate,
  so the cost of each grade of assurance stays a published number
  rather than a promise.
\end{itemize}

